% Options for packages loaded elsewhere
\PassOptionsToPackage{unicode}{hyperref}
\PassOptionsToPackage{hyphens}{url}
\documentclass[
  12pt,
]{ctexart}
\usepackage{xcolor}
\usepackage{amsmath,amssymb}
\setcounter{secnumdepth}{-\maxdimen} % remove section numbering
\usepackage{iftex}
\ifPDFTeX
  \usepackage[T1]{fontenc}
  \usepackage[utf8]{inputenc}
  \usepackage{textcomp} % provide euro and other symbols
\else % if luatex or xetex
  \usepackage{unicode-math} % this also loads fontspec
  \defaultfontfeatures{Scale=MatchLowercase}
  \defaultfontfeatures[\rmfamily]{Ligatures=TeX,Scale=1}
\fi
\usepackage{lmodern}
\ifPDFTeX\else
  % xetex/luatex font selection
\fi
% Use upquote if available, for straight quotes in verbatim environments
\IfFileExists{upquote.sty}{\usepackage{upquote}}{}
\IfFileExists{microtype.sty}{% use microtype if available
  \usepackage[]{microtype}
  \UseMicrotypeSet[protrusion]{basicmath} % disable protrusion for tt fonts
}{}
\makeatletter
\@ifundefined{KOMAClassName}{% if non-KOMA class
  \IfFileExists{parskip.sty}{%
    \usepackage{parskip}
  }{% else
    \setlength{\parindent}{0pt}
    \setlength{\parskip}{6pt plus 2pt minus 1pt}}
}{% if KOMA class
  \KOMAoptions{parskip=half}}
\makeatother
\usepackage{longtable,booktabs,array}
\usepackage{calc} % for calculating minipage widths
% Correct order of tables after \paragraph or \subparagraph
\usepackage{etoolbox}
\makeatletter
\patchcmd\longtable{\par}{\if@noskipsec\mbox{}\fi\par}{}{}
\makeatother
% Allow footnotes in longtable head/foot
\IfFileExists{footnotehyper.sty}{\usepackage{footnotehyper}}{\usepackage{footnote}}
\makesavenoteenv{longtable}
\setlength{\emergencystretch}{3em} % prevent overfull lines
\providecommand{\tightlist}{%
  \setlength{\itemsep}{0pt}\setlength{\parskip}{0pt}}
\usepackage{bookmark}
\IfFileExists{xurl.sty}{\usepackage{xurl}}{} % add URL line breaks if available
\urlstyle{same}
\hypersetup{
  hidelinks,
  pdfcreator={LaTeX via pandoc}}

\author{SCX}
\date{2026}

\begin{document}

\section{SCX/EGP 竞争分析报告：MoE
架构、蒸馏与专家治理}\label{scxegp-ux7adeux4e89ux5206ux6790ux62a5ux544amoe-ux67b6ux6784ux84b8ux998fux4e0eux4e13ux5bb6ux6cbbux7406}

\begin{quote}
生成时间：2026-06-25

覆盖范围：2023-2026 年最新文献，涵盖 Mixture-of-Experts (MoE)、Knowledge
Distillation (KD)、模型合并、MLIP gauge 对齐，以及专家治理（Expert
Governance）四大方向。

目标：评估每一项工作与 SCX/EGP 框架的关系，定位 SCX/EGP
的差异化优势及未被占领的技术空白。
\end{quote}

\begin{center}\rule{0.5\linewidth}{0.5pt}\end{center}

\subsection{目录}\label{ux76eeux5f55}

\begin{enumerate}
\def\labelenumi{\arabic{enumi}.}
\tightlist
\item
  \hyperref[a-mixture-of-experts-moe-in-mlip]{A. Mixture of Experts
  (MoE) in MLIP}
\item
  \hyperref[b-knowledge-distillation-for-mlip]{B. Knowledge Distillation
  for MLIP}
\item
  \hyperref[c-expert-merging--gauge--energy-alignment]{C. Expert Merging
  / Gauge / Energy Alignment}
\item
  \hyperref[d-expert-compilation--governancescx-ux6838ux5fc3ux521bux65b0ux533a]{D.
  Expert Compilation / Governance（SCX 核心创新区）}
\item
  \hyperref[e-ux7adeux4e89ux683cux5c40ux603bux7ed3ux4e0eux6218ux7565ux5efaux8bae]{E.
  竞争格局总结与战略建议}
\end{enumerate}

\begin{center}\rule{0.5\linewidth}{0.5pt}\end{center}

\subsection{A. Mixture of Experts (MoE) in
MLIP}\label{a.-mixture-of-experts-moe-in-mlip}

\subsubsection{A1. DeepMD MoE/MoLE --- ``Scaling MLIP with Mixtures of
Experts''}\label{a1.-deepmd-moemole-scaling-mlip-with-mixtures-of-experts}

\textbf{论文：} Liu, Zhang, Peng, E, Zhang, Wang. ``Scaling Machine
Learning Interatomic Potentials with Mixtures of Experts.''
arXiv:2603.07977 (2026).

\textbf{他们解决了什么？} 将 MoE 架构系统性地引入 MLIP 领域。此前 MoE 在
NLP 和 CV 中已取得巨大成功，但面向原子间势函数的 MoE
架构设计（路由策略、专家结构、稀疏激活）缺乏系统研究。

\textbf{核心方法：} -
提出了两种架构变体：\textbf{MoLE（Mixture-of-Linear-Experts）} 和
\textbf{MoE（非线性 MLP 专家）}。 -
设计三种路由策略对比：\textbf{逐元素路由}（element-wise routing）
vs.~构型级路由（configuration-level routing） vs.~全局 MoE 路由。 -
引入\textbf{共享专家（shared expert）} 机制增强泛化。 - 经由 DeepMD-kit
框架实现，基于 DP 模型架构扩展。

\textbf{关键发现：} - \textbf{逐元素路由（E(3) equivariant
routing）显著优于构型级路由}，全局 MoE 路由常导致数值不稳定。 - 非线性
MoE + 共享专家组合效果最优，超越 MoLE。 -
路由模式呈现\textbf{化学可解释的专家分工}：不同专家倾向于处理特定元素族（与周期表趋势一致）。
- 在 OMol25、OMat24、OC20M 基准上达到 SOTA。

\textbf{与 SCX/EGP 的关系：} - \textbf{强相关。} DeepMD MoE 是目前与
SCX/EGP 最直接的竞争方法。 - 二者都使用 MoE
思想来提升势函数精度。区别在于：DeepMD MoE 在\textbf{训练时}构建 MoE
架构（端到端训练一个模型），而 SCX/EGP
是\textbf{后训练}方式------将多个独立训练的专家模型编译为联合系统。 -
DeepMD MoE 的专家是在训练中自动形成的（emergent expertise），而 SCX/EGP
的专家是独立训练的（pre-defined expertise）。

\textbf{SCX 相对优势：} - \textbf{灵活性：} SCX/EGP
可以组合任意来源的独立训练模型（不同架构、不同训练集、不同超参数），而
DeepMD MoE 所有专家必须在同一框架下端到端训练。 - \textbf{可扩展性：}
EGP 的专家可以是异构的（MACE + NEP + SevenNet），DeepMD MoE
的所有专家必须是同构的 DP 架构。 - \textbf{增量更新：} EGP
可以随时加入新模型而无需重训，DeepMD MoE 需要整体重训。

\textbf{空白填补机会：} - DeepMD MoE
未讨论\textbf{专家选择的可解释性}------SCX/EGP
的输入条件路由可以明确解释''对给定结构，为何选择某专家''。 - DeepMD MoE
没有蒸馏流程------SCX/EGP 的蒸馏环节可与 MoE 结合形成''MoE 预编译 +
蒸馏''流水线。

\begin{center}\rule{0.5\linewidth}{0.5pt}\end{center}

\subsubsection{A2. Allegro/NequIP MoE --- Spatial Partition +
Co-training
Agreement}\label{a2.-allegronequip-moe-spatial-partition-co-training-agreement}

\textbf{论文：} Nascimento et al.~``Mixture of Experts Framework in
Machine Learning Interatomic Potentials for Atomistic Simulations.''
arXiv:2604.26143 (2026).

\textbf{他们解决了什么？} 解决了\textbf{空间分区 MoE}
中的关键问题：不同复杂度模型在界面处的力学失配（artificial stress
fields），以及高频模型在局部区域带来的计算瓶颈。

\textbf{核心方法：} - 基于 E(3)-equivariant \textbf{Allegro}
架构构建多保真 MoE 框架。 - \textbf{空间分区（Spatial Partitioning）：}
将模拟域分为化学复杂区域（如催化反应界面）和简单区域（如体相晶格），分配不同保真度的模型。
- \textbf{协同训练（Co-training）：}
损失函数包含\textbf{一致性约束（agreement
constraints）}------对共享体相环境的逐原子能量和力差异施加惩罚，迫使不同专家学习一致的体相物理描述。
- 验证体系：Pt+CO 催化系统。

\textbf{结果：} - 精确能量守恒（exact energy conservation maintained） -
体相力学响应（EOS、体模量）在专家间对齐 - 预测精度与全高保真模拟相当 -
\textbf{计算速度提升超过 2 倍}

\textbf{与 SCX/EGP 的关系：} - \textbf{互补关系。}
该工作解决的是\textbf{空间分区}问题（不同区域用不同模型），SCX/EGP
解决的是\textbf{输入条件路由}问题（不同条件用不同模型）。 -
二者的共识在于：\textbf{单一模型不足以覆盖所有场景}，需要多个专家协作。
- 该工作的''一致性约束''思想与 SCX/EGP 的''Gauge
对称化''有关联------都试图使多个模型在重叠区域保持一致预测。

\textbf{SCX 相对优势：} - SCX/EGP
不依赖空间分区假设，因此适用于\textbf{非局域问题}（如长程相互作用、全局结构变化），而空间分区
MoE 仅适用于可划分的体系。 - SCX/EGP
的专家数量不受空间维度限制，可扩展到任意多专家。

\textbf{空白填补机会：} -
该工作的''一致性约束''需要通过额外训练来实现；SCX/EGP
可以在\textbf{不重训}的情况下实现专家间一致性（通过校准层的后处理）。 -
空间分区与输入条件路由可以结合------例如在界面上同时使用两种策略。

\begin{center}\rule{0.5\linewidth}{0.5pt}\end{center}

\subsubsection{A3. 通用 MoE --- Shazeer et al.~2017, Switch
Transformer}\label{a3.-ux901aux7528-moe-shazeer-et-al.-2017-switch-transformer}

\textbf{论文：} - Shazeer et al.~``Outrageously Large Neural Networks:
The Sparsely-Gated Mixture-of-Experts Layer.'' ICLR 2017. - Fedus, Zoph
\& Shazeer. ``Switch Transformers: Scaling to Trillion Parameter Models
with Simple and Efficient Sparsity.'' JMLR 2022.

\textbf{他们解决了什么？} Shazeer 2017 提出了\textbf{稀疏门控 MoE
层}的核心架构：通过一个可学习的门控网络将每个输入路由到 top-k
个专家（子网络），实现''总参数量巨大但计算量恒定''的效果。Switch
Transformer 将其简化为 top-1 路由，显著降低了通信和计算开销。

\textbf{核心方法：} - \textbf{Noisy Top-K Gating：} 在门控 logits
上加噪声后取 top-k，实现稀疏激活。 - \textbf{辅助负载均衡损失（Switch
Transformer）：}
通过不同化的负载均衡损失避免''富者愈富''（少数专家支配大部分输入）问题。
- \textbf{容量因子（Capacity Factor）：} 控制每个专家处理的 token
数上限。 - \textbf{训练稳定技巧：} 选择性精度（router 用
float32）、更小的初始化、专家 dropout。

\textbf{影响：} 直接催生了 Mixtral 8x7B、DeepSeek-V2/V3、Llama 4 等现代
MoE LLM。

\textbf{与 SCX/EGP 的关系：} - \textbf{理论基础关系。} 这些工作奠定了
MoE 路由、负载均衡的理论基础。SCX/EGP
的''输入条件路由''本质上是同一问题在势函数领域的应用，但面临不同挑战： -
NLP MoE 的输入是离散 token，路由空间相对有限 - SCX/EGP
的输入是连续原子构型，路由空间是高维连续流形 - NLP MoE 的专家是同构 FFN
层；SCX/EGP 的专家可以是异构的完整势函数

\textbf{SCX 相对优势：} - SCX/EGP
的核心理念------\textbf{``训练时不耦合，推理时再融合''}------避免了通用
MoE 的两个致命问题： 1. \textbf{训练不稳定：} 通用 MoE
的路由器与专家耦合训练，常导致专家坍塌（collapse）或不稳定。 2.
\textbf{负载均衡的开销：} 需要复杂的辅助损失函数。 - SCX/EGP
中每个专家可独立训练和验证，质量可控。

\textbf{空白填补机会：} - 通用 MoE
文献中丰富的理论成果（负载均衡、容量规划、路由策略）可以迁移到 SCX/EGP
框架中。 - 思考：能否借鉴 NLP MoE 的''共享专家''（shared expert）思想到
SCX/EGP 中？

\begin{center}\rule{0.5\linewidth}{0.5pt}\end{center}

\subsection{B. Knowledge Distillation for
MLIP}\label{b.-knowledge-distillation-for-mlip}

\subsubsection{B1. MACE → NEP 蒸馏}\label{b1.-mace-nep-ux84b8ux998f}

\textbf{论文：} ``Constructing machine learning interatomic potentials
with minimum amount of ab initio data.'' npj Computational Materials,
s41524-026-02023-y (2026).

\textbf{他们解决了什么？} 用最少量的 DFT 数据（\textasciitilde200
个构型）构建高精度 MLIP，无需昂贵的主动学习循环（通常需要数千个 DFT
计算）。

\textbf{核心方法：} - \textbf{两阶段教师-学生蒸馏：} - Stage
1（教师）：从预训练通用模型 \textbf{MACE-MP-0} 开始，在
\textasciitilde200 个 DFT 构型上微调。 - Stage 2（学生）：将微调后的
MACE 教师蒸馏到轻量 \textbf{NEP} 模型。 -
验证体系：固态电解质（LGPS、LATP、LYC）。

\textbf{优势：} 单次（single-shot）工作流，无需主动学习循环。NEP
的推理速度远快于 MACE，适合大规模 MD 模拟。

\textbf{与 SCX/EGP 的关系：} - \textbf{强相关。} 该工作展示了''昂贵教师
→ 廉价学生''的标准蒸馏流水线。SCX/EGP 的蒸馏环节本质上也是''复杂专家 →
统一学生''，但 SCX/EGP 的关键差异在于： 1. \textbf{蒸馏前需规范化（Gauge
对齐）：} SCX/EGP 的 Gauge
对称化确保专家间能量零点一致，而该工作假设单一教师没有问题。 2.
\textbf{多专家蒸馏：} SCX/EGP
面临的不是单教师→单学生，而是多教师→单学生。

\textbf{SCX 相对优势：} - SCX/EGP
处理多专家情况下的''冲突预测''问题，该工作未涉及。 - SCX/EGP
可以蒸馏来自不同架构类型（MACE、NequIP、SevenNet
等）的教师，而该工作只有单一架构。

\begin{center}\rule{0.5\linewidth}{0.5pt}\end{center}

\subsubsection{B2. Ensemble Knowledge Distillation (EKD) for
MLIP}\label{b2.-ensemble-knowledge-distillation-ekd-for-mlip}

\textbf{论文：} Matin et al.~``Ensemble Knowledge Distillation for
Machine Learning Interatomic Potentials.'' arXiv:2503.14293 (2025).

\textbf{他们解决了什么？} 解决能量数据集缺少力标签的问题。在只有 QC
能量的数据集上，通过集成多个教师模型生成伪力，提高学生模型的精度。

\textbf{核心方法：} 1. 在 QC 能量上训练 \textbf{多个教师
MLIP}（多种初始化/架构）。 2.
用集成教师生成\textbf{平均伪力}（ensemble-averaged forces）作为软标签。
3. 学生模型同时学习 \textbf{QC 能量 + 集成平均力}。 4. 验证于 ANI-1ccx
数据集（耦合簇级别能量），在 COMP6 基准上达到 SOTA。

\textbf{与 SCX/EGP 的关系：} - \textbf{弱相关但有启发性。}
该工作关注的是''如何生成更好的软标签''，SCX/EGP
关注的是''如何组合多个独立训练的专家''。EKD
自然地使用了集成思路，但目的是蒸馏而非专家协作。

\textbf{SCX 相对优势：} - EKD 的所有教师需要共享同一训练数据；SCX/EGP
的专家可以在不同数据上独立训练。 - EKD
的学生架构不能是教师集的一部分；SCX/EGP
的''学生''本质上是编译器化的专家系统，保留了专家多样性。

\begin{center}\rule{0.5\linewidth}{0.5pt}\end{center}

\subsubsection{B3. Foundation Model
Distillation}\label{b3.-foundation-model-distillation}

\textbf{论文：} Gardner et al.~``Distillation of atomistic foundation
models across architectures and chemical domains.'' arXiv:2506.10956
(2025).

\textbf{他们解决了什么？} 证明合成数据蒸馏（distillation via synthetic
data）可以便宜地将原子级基础模型的知识迁移到不同架构。

\textbf{核心方法：} - 教师：基础模型（large universal MLIP） -
生成合成数据，标记后训练轻量学生 - 演示 \textgreater10x（GNN→GNN）和
\textgreater100x（GNN→ACE）加速 -
验证体系：液态水、极端条件下的氢、多孔二氧化硅、杂化钙钛矿、有机反应

\textbf{与 SCX/EGP 的关系：} - \textbf{中等相关。}
这是''基础模型→专门模型''的知识迁移，SCX/EGP
是''多个专门模型→编译器化系统''的专家协作。二者互补------基础模型可以作为
SCX/EGP 的专家候选池之一。

\textbf{空白填补机会：} - 该工作的''跨架构蒸馏''可以借鉴到 SCX/EGP
的蒸馏环节（如何从异构专家蒸馏到统一学生）。 -
目前没有工作探讨''如何从多个基础模型中蒸馏''------这正是 SCX/EGP
的用武之地。

\begin{center}\rule{0.5\linewidth}{0.5pt}\end{center}

\subsubsection{B4. Teacher-Student Training for MLIPs (RSC
2025)}\label{b4.-teacher-student-training-for-mlips-rsc-2025}

\textbf{论文：} Matin et al.~``Teacher-student training improves the
accuracy and efficiency of machine learning interatomic potentials.''
RSC Digital Discovery, 2025.

\textbf{他们解决了什么？}
证明教师-学生训练不仅能加速推理，还能\textbf{提高学生模型的准确率}（甚至超过教师）。

\textbf{核心方法：} - 架构：\textbf{HIPNN}（Hierarchically Interacting
Particle Neural Network） -
损失函数增加一项：匹配学生与教师的逐原子能量预测 -
辅助标签：教师网络隐藏层的原子能量输出 - 实验：学生模型实现
\textbf{\textgreater2 倍速度}、\textbf{\textless50\%
内存}，且精度\textbf{超越}教师 - \textbf{帕累托主导}（Pareto
dominance）：在每个成本点上都优于对照模型

\textbf{与 SCX/EGP 的关系：} - \textbf{中等相关。}
关键启示：\textbf{蒸馏可以同时提高精度和效率}。SCX/EGP
的蒸馏环节应追求类似效果，而不仅仅是为了压缩。

\textbf{空白填补机会：} - 该工作仅使用单一教师。SCX/EGP
有机会展示''多教师蒸馏''不仅可以压缩模型，还可以通过集成效应超越任何单个专家------即
\textbf{EGP 的协同效应（synergy）}。

\begin{center}\rule{0.5\linewidth}{0.5pt}\end{center}

\subsubsection{B5.
其他重要蒸馏工作}\label{b5.-ux5176ux4ed6ux91cdux8981ux84b8ux998fux5de5ux4f5c}

\begin{longtable}[]{@{}
  >{\raggedright\arraybackslash}p{(\linewidth - 6\tabcolsep) * \real{0.1765}}
  >{\raggedright\arraybackslash}p{(\linewidth - 6\tabcolsep) * \real{0.1765}}
  >{\raggedright\arraybackslash}p{(\linewidth - 6\tabcolsep) * \real{0.1765}}
  >{\raggedright\arraybackslash}p{(\linewidth - 6\tabcolsep) * \real{0.4706}}@{}}
\toprule\noalign{}
\begin{minipage}[b]{\linewidth}\raggedright
工作
\end{minipage} & \begin{minipage}[b]{\linewidth}\raggedright
年份
\end{minipage} & \begin{minipage}[b]{\linewidth}\raggedright
贡献
\end{minipage} & \begin{minipage}[b]{\linewidth}\raggedright
与 SCX/EGP 关系
\end{minipage} \\
\midrule\noalign{}
\endhead
\bottomrule\noalign{}
\endlastfoot
\textbf{SevenNet-Nano} (arXiv:2604.10887) & 2026 & 从 SevenNet-Omni
蒸馏到 Nano，\textgreater10 倍加速，极端条件验证 &
补偿性：单一模型蒸馏 \\
\textbf{PFD 框架} (Phys. Rev.~Materials, 2025) & 2025 &
Pre-train→Fine-tune→Distill 流水线，\textasciitilde100 DFT
帧训练\textasciitilde100x加速学生 & 流程互补：PFD 是蒸馏前的准备流程 \\
\textbf{ARK Distillation} (npj Comput. Mater., 2026) & 2026 &
角度关系知识蒸馏 + 对比学习，11.9x 催化剂筛选加速 &
方法创新：蒸馏中的结构保持 \\
\end{longtable}

\textbf{ARK Distillation 特别值得关注：} - 论文：Lim et al.~``Angular
relational knowledge distillation of machine learning interatomic
potentials for scalable catalyst exploration.'' npj Computational
Materials, 2026. - 核心创新：用\textbf{角度关系向量（angular relational
vectors）} 和\textbf{对比学习}保持教师 PES 的几何结构。 - 高通量筛选：58
万结构筛选仅用 11.6 GPU-hours（对照：59.4 CPU-years）。 - 启示：SCX/EGP
的蒸馏环节可以借鉴这种''结构保持''的思想，确保编译后的学生模型不丢失专家对
PES 几何的精细刻画。

\textbf{PFD 框架特别值得关注：} - 论文：Wang, Gao et al.~``Pre-training,
fine-tuning, and distillation (PFD): Automatically generating machine
learning force fields from universal models.'' Phys. Rev.~Materials,
2025. - 核心：DPA-2 → 微调 → DeePMD 学生，仅 \textasciitilde100 帧 DFT
数据。 - SCX/EGP 可以与 PFD 结合：PFD 作为''专家生成''环节，SCX/EGP
作为''专家编译''环节。

\begin{center}\rule{0.5\linewidth}{0.5pt}\end{center}

\subsection{C. Expert Merging / Gauge / Energy
Alignment}\label{c.-expert-merging-gauge-energy-alignment}

\subsubsection{C1. Model Merging (Model Soup, Git Re-Basin,
线性模式连接)}\label{c1.-model-merging-model-soup-git-re-basin-ux7ebfux6027ux6a21ux5f0fux8fdeux63a5}

\textbf{核心思想：}
将多个独立训练的模型权重直接平均（或经过排列对齐后平均）获得合并模型，无需再训练。核心假设：不同模型的解在权重空间中通过线性路径连接且损失面平坦。

\textbf{关键工作：}

\begin{longtable}[]{@{}
  >{\raggedright\arraybackslash}p{(\linewidth - 4\tabcolsep) * \real{0.3333}}
  >{\raggedright\arraybackslash}p{(\linewidth - 4\tabcolsep) * \real{0.3333}}
  >{\raggedright\arraybackslash}p{(\linewidth - 4\tabcolsep) * \real{0.3333}}@{}}
\toprule\noalign{}
\begin{minipage}[b]{\linewidth}\raggedright
工作
\end{minipage} & \begin{minipage}[b]{\linewidth}\raggedright
年份
\end{minipage} & \begin{minipage}[b]{\linewidth}\raggedright
贡献
\end{minipage} \\
\midrule\noalign{}
\endhead
\bottomrule\noalign{}
\endlastfoot
\textbf{Model Soups} (Wortsman et al., ICML 2022) & 2022 &
均匀/贪心平均微调模型权重，提升精度且不增推理时间 \\
\textbf{Git Re-Basin} (Ainsworth et al., ICLR 2023) & 2023 &
通过排列对齐（permutation alignment）使不同模型处于同一
basin，再平均权重 \\
\textbf{ZipIt!} (ICLR 2024) & 2024 &
可以合并来自不同任务、不同架构的模型（以特征融合代替简单平均） \\
\textbf{AME} (Lee \& Ngo, NeurIPS 2025) & 2025 & 将模型 soup
重新解释为''摊销模型集成（Amortized Model
Ensembling）``，使用伪梯度进行无数据元优化 \\
\textbf{Soup-of-Experts} (ICML 2025) & 2025 &
参数级专家线性组合，测试时可即时实例化任意域权重组合的模型 \\
\textbf{LLM Souping} (Meta, 2025) & 2025 & 加权平均多个微调
LLM，引入类别感知专家选择 SoCE \\
\end{longtable}

\textbf{与 SCX/EGP 的关系：} - \textbf{概念近但方法不同。} 模型合并和
SCX/EGP 都追求''结合多个模型的知识''。但： -
模型合并在\textbf{权重空间}操作（weight-space
averaging），要求模型架构相同。 - SCX/EGP
在\textbf{预测空间}操作（prediction-space ensembling），可处理异构架构。
- 模型合并是一个\textbf{融合}过程（多个模型融合为一个）；SCX/EGP
是一个\textbf{编译}过程（多个模型保持独立但通过路由系统协调）。

\textbf{SCX 相对优势：} - \textbf{架构无关性：}
模型合并要求同构架构，SCX/EGP 可处理 MACE + NEP + SevenNet 异构组合。 -
\textbf{可竞争性：} 模型合并建立于''模型在权重空间中相容''的假设上，这在
MLIP 中是否成立尚待验证。SCX/EGP 不需要这种假设。 - \textbf{可解释性：}
SCX/EGP 保留了每个专家的身份，可以追溯哪些专家对哪些输入负责。

\textbf{空白填补机会：} - 目前\textbf{没有工作}将模型合并方法系统应用于
MLIP 领域。 - 可以尝试''MLIP
在权重空间中的线性模式连通性''研究------这是一个开放问题。

\begin{center}\rule{0.5\linewidth}{0.5pt}\end{center}

\subsubsection{C2. MLIP Gauge Freedom 与 Species Energy
Shift}\label{c2.-mlip-gauge-freedom-ux4e0e-species-energy-shift}

\textbf{论文：} ``Gauge dependence of total energies and
cross-functional transfer learning.'' arXiv:2504.05565 (2025).

\textbf{他们解决了什么？} 明确指出了 MLIP 中的 \textbf{Gauge 自由（gauge
freedom）} 问题：DFT
总能依赖于真空能级的标度选择，不同泛函的总能量基准不同。这个''Gauge''的选择影响
MLIP 的迁移学习效果。

\textbf{核心方法：} - 分析多种迁移学习策略，对比''是否调整
AtomRef（species energy shift）``的影响。 - \textbf{关键发现：}
跨功能迁移学习时，\textbf{先 refit AtomRef 再训练 GNN}
是最优策略。仅重训练 AtomRef 就将 Pearson 相关系数从 0.09 提升到 0.93。
- 具体结果：Method 4（refit AtomRef + train GNN）能量 MAE = 17
meV/atom，而直接在源功能基上训练 GNN 是 26 meV/atom。

\textbf{与 SCX/EGP 的关系：} - \textbf{核心相关。} Gauge freedom 正是
SCX/EGP 中 \textbf{Gauge 对称化}环节要解决的问题。不同 MLIP
在相同结构上的总能量预测存在系统性能量偏移（例如，OUTCAR 中的''energy
without entropy''和 MACE 的预测相差数 eV）。 - SCX/EGP 的 Gauge
层（GAU）作用就是消除这种偏移，使不同专家的能量预测在联合前可比较。

\textbf{SCX 相对优势：} - 该工作需要每个专家都有自己的 AtomRef
且需要在数据集上显式计算。SCX/EGP 的 Gauge
层可以在不接触原始训练数据的情况下\textbf{在线校准}------只需在重叠构型上比较预测值。
- SCX/EGP 的 Gauge
对称化不仅处理能量偏移，还处理\textbf{力和应力的对齐}（通过 DOF 变换）。

\textbf{空白填补机会：} - 该工作中指出''formation energy''方法是
Gauge-invariant 的替代方案。SCX/EGP 可以探索：是否在 formation energy
空间而非 total energy 空间进行专家编译？ -
目前没有工作系统地探讨\textbf{多专家之间的 Gauge 不一致性}------SCX/EGP
是第一个系统解决该问题的框架。 - 该工作只跨功能基（GGA/GGA+U →
r²SCAN）迁移。SCX/EGP 可以扩展到跨势函数的 Gauge 对齐（SPME → MACE →
NEP）。

\begin{center}\rule{0.5\linewidth}{0.5pt}\end{center}

\subsubsection{C3. Species Energy Shift
标准做法}\label{c3.-species-energy-shift-ux6807ux51c6ux505aux6cd5}

当前主流通用 MLIP 中的能量偏移处理：

\begin{longtable}[]{@{}lll@{}}
\toprule\noalign{}
模型 & 方法 & 说明 \\
\midrule\noalign{}
\endhead
\bottomrule\noalign{}
\endlastfoot
CHGNet & AtomRef + GNN & 元素级线性偏移 + 图神经网络残差 \\
M3GNet & AtomRef + GNN & 同上 \\
NequIP & AtomRef (subtract per-atom energy) & 训练前减去参考能量 \\
CACE & AtomRef + GNN & Composible Atomic Cluster Expansion \\
MACE-MP-0 & 平均原子能偏移 & 基于 MP 数据集计算平均偏移 \\
SevenNet & 隐式 offset & 模型自动学习偏移 \\
\end{longtable}

所有这些方法都是\textbf{单模型校准}------将单个模型的能量预测对齐到某个参考。SCX/EGP
的 Gauge 对称化是第一个处理\textbf{跨模型能量对齐}的方法。

\begin{center}\rule{0.5\linewidth}{0.5pt}\end{center}

\subsection{D. Expert Compilation / Governance（SCX
核心创新区）}\label{d.-expert-compilation-governancescx-ux6838ux5fc3ux521bux65b0ux533a}

这是 SCX/EGP
最核心的差异化领域。搜索结果显示，该方向目前\textbf{基本没有竞争工作}。以下是搜索覆盖的几个子方向：

\subsubsection{D1. ``蒸馏前的专家规范化''（Expert Normalization Before
Distillation）}\label{d1.-ux84b8ux998fux524dux7684ux4e13ux5bb6ux89c4ux8303ux5316expert-normalization-before-distillation}

\textbf{搜索结论：没有专门的工作。}

\begin{itemize}
\tightlist
\item
  现有蒸馏工作（MACE→NEP, SevenNet-Nano, PFD, ARK, EKD
  等）都假设单一教师或同构集成教师。
\item
  当有多教师时，现有工作简单地平均伪力（如
  EKD）或者忽略多教师之间的能量不一致性。
\item
  \textbf{没有工作系统地解决''蒸馏前必须对齐专家的能量标度''这一问题。}
\item
  SCX/EGP 的 Gauge → Route → Compile
  流水线是唯一明确区分''专家规范化''与''专家合并''的框架。
\end{itemize}

\textbf{SCX 的机会：} - 这是 SCX/EGP
最明确的技术壁垒。如果有其他工作开始研究''多教师蒸馏前的能量对齐''，SCX
需要及时跟进和引用。

\begin{center}\rule{0.5\linewidth}{0.5pt}\end{center}

\subsubsection{D2. Expert Domain
Certification}\label{d2.-expert-domain-certification}

\textbf{搜索结论：没有适用于 MLIP 的专家域认证工作。}

\begin{itemize}
\tightlist
\item
  通用 AI 领域有一些关于''AI 伦理专家认证''（AIGP
  等）的社会学/管理学研究，但这些与 MLIP 专家认证无关。
\item
  ``What Makes An Expert?'' (arXiv:2411.00179, 2024) 是一篇综述，讨论了
  ML 文献中''专家''定义模糊的问题，但聚焦在数据标注场景。
\item
  \textbf{没有工作定义''MLIP 专家在哪些输入域上可信''这一概念。}
\item
  SCX/EGP 的 ECR
  层（代表认证）所做的正是：评估每个专家在给定输入条件下的置信度。这是
  SCX 的原创贡献。
\end{itemize}

\textbf{SCX 的机会：} - 可以考虑发表专门的方法论文定义''MLIP Expert
Domain Certification''的形式化框架。 - 可以借鉴 conformal prediction
的思想，为每个专家在认证域内提供保证。

\begin{center}\rule{0.5\linewidth}{0.5pt}\end{center}

\subsubsection{D3. Multi-Expert Conflict
Arbitration（多专家冲突仲裁）}\label{d3.-multi-expert-conflict-arbitrationux591aux4e13ux5bb6ux51b2ux7a81ux4ef2ux88c1}

\textbf{搜索结论：在通用 ML 领域有探索，但在 MLIP 领域为空白。}

通用 ML 领域的相关工作：

\begin{longtable}[]{@{}
  >{\raggedright\arraybackslash}p{(\linewidth - 4\tabcolsep) * \real{0.3333}}
  >{\raggedright\arraybackslash}p{(\linewidth - 4\tabcolsep) * \real{0.3333}}
  >{\raggedright\arraybackslash}p{(\linewidth - 4\tabcolsep) * \real{0.3333}}@{}}
\toprule\noalign{}
\begin{minipage}[b]{\linewidth}\raggedright
工作
\end{minipage} & \begin{minipage}[b]{\linewidth}\raggedright
类型
\end{minipage} & \begin{minipage}[b]{\linewidth}\raggedright
方法
\end{minipage} \\
\midrule\noalign{}
\endhead
\bottomrule\noalign{}
\endlastfoot
\textbf{King Solomon Protocol} (2026) & 人类监督仲裁 &
正式的仲裁框架，双助手架构，结构化人类仲裁层 \\
\textbf{HCDF} (2026) & 层次化认知 &
知识获取→显式风险表示→结构仲裁三层架构 \\
\textbf{Cognitive Hives} (2026) & 分布式认知 &
专门的冲突解决层，识别/仲裁/协调争议 \\
\textbf{LLMediator / Mediator} (2025) & LLM 调解 &
基于不确定性路由，低冲突层平均 + 高冲突层 MoE \\
\textbf{AEGIS-Nexus} (2026) & 伦理仲裁 &
元司法仲裁代理，平衡影响力、可靠性和严重性 \\
\textbf{Multi-Expert Fusion Networks} survey & 综述 &
软门控、稀疏路由、top-k 仲裁等分类 \\
\end{longtable}

\textbf{MLIP 领域：没有工作探讨多势函数之间的预测冲突与仲裁。}

\begin{itemize}
\tightlist
\item
  现有的''势函数混合''方法（如 ML-MIX, 空间分区
  MoE）都回避了''多个模型在同一输入上给出不同预测时该怎么办''的问题。
\item
  SCX/EGP 的 Route
  层（RTE）本质上是一个仲裁器------它不仅选择专家，还整合专家的预测。
\end{itemize}

\textbf{SCX 的机会：} - 可以将通用 ML 仲裁的成果（如 HCDF
的层次化解耦、LLMediator 的不确定性路由）引入 MLIP 领域。 -
可以发表案例研究，展示 EGP 仲裁策略对预测质量的影响。

\begin{center}\rule{0.5\linewidth}{0.5pt}\end{center}

\subsubsection{D4.
多专家性能的输入条件依赖}\label{d4.-ux591aux4e13ux5bb6ux6027ux80fdux7684ux8f93ux5165ux6761ux4ef6ux4f9dux8d56}

\textbf{搜索结论：没有系统性的工作。}

\begin{itemize}
\tightlist
\item
  MoE
  文献通常关注\textbf{平均性能}，而不是''对特定类型的输入，哪个专家更好''。
\item
  Active learning 和 ALIP
  等数据选择方法隐含了''某些区域某些模型更好''但未形式化。
\item
  DeepMD MoE
  的逐元素路由展示了''不同的元素的构型应由不同专家处理''------这是一个好的开端。
\item
  \textbf{没有工作系统地构建''专家性能配置文件''（expert performance
  profile over input conditions）。}
\end{itemize}

\textbf{SCX 的机会：} - ECR
层本质上就是一个专家性能配置文件的实时评估器。 -
可以发布一个专家性能配置文件的自动构建工具。

\begin{center}\rule{0.5\linewidth}{0.5pt}\end{center}

\subsection{E.
竞争格局总结与战略建议}\label{e.-ux7adeux4e89ux683cux5c40ux603bux7ed3ux4e0eux6218ux7565ux5efaux8bae}

\subsubsection{E1.
创新维度对比矩阵}\label{e1.-ux521bux65b0ux7ef4ux5ea6ux5bf9ux6bd4ux77e9ux9635}

\begin{longtable}[]{@{}
  >{\raggedright\arraybackslash}p{(\linewidth - 10\tabcolsep) * \real{0.2857}}
  >{\centering\arraybackslash}p{(\linewidth - 10\tabcolsep) * \real{0.1429}}
  >{\centering\arraybackslash}p{(\linewidth - 10\tabcolsep) * \real{0.1429}}
  >{\centering\arraybackslash}p{(\linewidth - 10\tabcolsep) * \real{0.1429}}
  >{\centering\arraybackslash}p{(\linewidth - 10\tabcolsep) * \real{0.1429}}
  >{\centering\arraybackslash}p{(\linewidth - 10\tabcolsep) * \real{0.1429}}@{}}
\toprule\noalign{}
\begin{minipage}[b]{\linewidth}\raggedright
竞争维度
\end{minipage} & \begin{minipage}[b]{\linewidth}\centering
DeepMD MoE (2603)
\end{minipage} & \begin{minipage}[b]{\linewidth}\centering
Allegro MoE (2604)
\end{minipage} & \begin{minipage}[b]{\linewidth}\centering
模型合并
\end{minipage} & \begin{minipage}[b]{\linewidth}\centering
MLIP 蒸馏
\end{minipage} & \begin{minipage}[b]{\linewidth}\centering
SCX/EGP
\end{minipage} \\
\midrule\noalign{}
\endhead
\bottomrule\noalign{}
\endlastfoot
训练时 MoE & Y & Y & N & N & N \\
后训练/编译 & N & N & Y & Y & \textbf{Y} \\
异构专家 & N & N & N & N & \textbf{Y} \\
Gauge 对齐 & N & Y(隐式) & N & N & \textbf{Y} \\
输入条件路由 & Y(元素级) & Y(空间) & N & N & \textbf{Y} \\
专家域认证 & N & N & N & N & \textbf{Y} \\
蒸馏前规范化 & N & N & N & N & \textbf{Y} \\
冲突仲裁 & N & N & N & N & \textbf{Y} \\
训练不可知 & N & N & Y & N & \textbf{Y} \\
\end{longtable}

\textbf{核心结论：SCX/EGP 在''后训练编译''这个细分领域没有直接竞争。}
每个现有工作都只在某些维度上与 SCX/EGP 有交集，但没有任何工作覆盖
SCX/EGP 的全部创新维度。

\subsubsection{E2.
竞争风险等级评估}\label{e2.-ux7adeux4e89ux98ceux9669ux7b49ux7ea7ux8bc4ux4f30}

\begin{longtable}[]{@{}
  >{\centering\arraybackslash}p{(\linewidth - 4\tabcolsep) * \real{0.2381}}
  >{\raggedright\arraybackslash}p{(\linewidth - 4\tabcolsep) * \real{0.4762}}
  >{\raggedright\arraybackslash}p{(\linewidth - 4\tabcolsep) * \real{0.2857}}@{}}
\toprule\noalign{}
\begin{minipage}[b]{\linewidth}\centering
风险等级
\end{minipage} & \begin{minipage}[b]{\linewidth}\raggedright
竞争方向
\end{minipage} & \begin{minipage}[b]{\linewidth}\raggedright
原因
\end{minipage} \\
\midrule\noalign{}
\endhead
\bottomrule\noalign{}
\endlastfoot
\textbf{高} & DeepMD MoE (2603.07977) & DeepMD
社区大、影响力强，如果他们将路由策略扩展到后训练场景，将是直接竞争 \\
\textbf{高} & Allegro MoE (2604.26143) &
如果他们将''空间分区''扩展到''输入条件分区''，会与 EGP 的路由层重叠 \\
\textbf{中} & ARK 蒸馏 (npj 2026) & 如果 ARK
扩展到多教师蒸馏并引入对齐步骤，将进入''蒸馏前规范化''领域 \\
\textbf{中} & PFD 流程 (PRM 2025) & 如果 PFD
扩展到''多模型联合蒸馏''，将接近 EGP 的 Compile 层 \\
\textbf{低} & 模型合并 (通用 ML) &
架构不可知方向不同，且合并需要同构架构 \\
\textbf{低} & Foundation 蒸馏 (2506.10956) &
目前聚焦单教师→单学生，进入多教师领域还需较大跨度 \\
\end{longtable}

\subsubsection{E3. SCX/EGP
必须主动占领的空白}\label{e3.-scxegp-ux5fc5ux987bux4e3bux52a8ux5360ux9886ux7684ux7a7aux767d}

按优先级排序：

\begin{enumerate}
\def\labelenumi{\arabic{enumi}.}
\tightlist
\item
  \textbf{蒸馏前的专家规范化（Prerequisite Normalization）}

  \begin{itemize}
  \tightlist
  \item
    这是 SCX/EGP 最独特、最差异化的贡献。
  \item
    必须通过论文、代码、演示案例让社区认识到：\textbf{跨专家能量对齐是多专家编译的必要前提}。
  \end{itemize}
\item
  \textbf{专家域认证（Expert Domain Certification）}

  \begin{itemize}
  \tightlist
  \item
    定义什么是''可信的专家''。
  \item
    建立基于 conformal prediction 或不确定性量化的认证框架。
  \item
    潜在应用：在提交 VASP
    任务前做专家域检查，避免在专家不确定的区域做不可靠预测。
  \end{itemize}
\item
  \textbf{跨架构专家编译演示（Cross-Architecture Compilation）}

  \begin{itemize}
  \tightlist
  \item
    展示 SCX/EGP 组合 MACE + NEP + SevenNet + NequIP 的能力。
  \item
    这是任何竞争对手都无法快速复制的技术壁垒------因为其他方法如模型合并、DeepMD
    MoE 都受限于架构同质性。
  \end{itemize}
\item
  \textbf{多专家蒸馏的协同效应（Synergistic Distillation）}

  \begin{itemize}
  \tightlist
  \item
    证明蒸馏后的 EGP 学生可以超越任何单个专家的精度。
  \item
    建立''EGP 蒸馏 \textgreater{} 单个专家微调 \textgreater{}
    直接蒸馏''的论证。
  \end{itemize}
\item
  \textbf{输入条件依赖的性能评估基准（Input-Conditioned Benchmark）}

  \begin{itemize}
  \tightlist
  \item
    发布在''专家更擅长什么''问题上的系统评估。
  \item
    可基于 NEB 路径、声子谱、弹性常数等多维属性。
  \item
    建立''哪个模型处理哪种类型的最佳''的参考标准。
  \end{itemize}
\end{enumerate}

\subsubsection{E4. 建议行动项}\label{e4.-ux5efaux8baeux884cux52a8ux9879}

\begin{longtable}[]{@{}
  >{\centering\arraybackslash}p{(\linewidth - 4\tabcolsep) * \real{0.2941}}
  >{\raggedright\arraybackslash}p{(\linewidth - 4\tabcolsep) * \real{0.3529}}
  >{\raggedright\arraybackslash}p{(\linewidth - 4\tabcolsep) * \real{0.3529}}@{}}
\toprule\noalign{}
\begin{minipage}[b]{\linewidth}\centering
优先级
\end{minipage} & \begin{minipage}[b]{\linewidth}\raggedright
行动
\end{minipage} & \begin{minipage}[b]{\linewidth}\raggedright
说明
\end{minipage} \\
\midrule\noalign{}
\endhead
\bottomrule\noalign{}
\endlastfoot
P0 & SCX 系统论文 & 在 arXiv 上发表 SCX/EGP 的方法论文，建立''Expert
Compilation''这一新方向的官方名称 \\
P0 & 开源 EGP 演示 & 基于现有的 VASP + MACE + NEP 等模型展示完整的
Gauge→Route→Compile 流水线 \\
P1 & Gauge 对称化论文 & 单独发表对 MLIP Gauge freedom
的量化分析和校准方法 \\
P1 & 输入条件路由 vs.~DeepMD MoE 对比基准 & 在 OMat24 等标准基准上与
DeepMD MoE 对比 \\
P2 & 专家域认证白皮书 & 定义 Expert Domain Certification 的形式化框架 \\
P2 & 冲突仲裁方法论论文 & 展示 EGP 的 Route 层如何处理多专家分歧 \\
\end{longtable}

\subsubsection{E5.
竞争方法背后的团队/组织}\label{e5.-ux7adeux4e89ux65b9ux6cd5ux80ccux540eux7684ux56e2ux961fux7ec4ux7ec7}

\begin{longtable}[]{@{}
  >{\raggedright\arraybackslash}p{(\linewidth - 4\tabcolsep) * \real{0.2857}}
  >{\raggedright\arraybackslash}p{(\linewidth - 4\tabcolsep) * \real{0.4762}}
  >{\centering\arraybackslash}p{(\linewidth - 4\tabcolsep) * \real{0.2381}}@{}}
\toprule\noalign{}
\begin{minipage}[b]{\linewidth}\raggedright
团队
\end{minipage} & \begin{minipage}[b]{\linewidth}\raggedright
相关论文
\end{minipage} & \begin{minipage}[b]{\linewidth}\centering
对 SCX 的竞争风险
\end{minipage} \\
\midrule\noalign{}
\endhead
\bottomrule\noalign{}
\endlastfoot
\textbf{DP Technology / AIS} (张林峰, 王涵等) & DeepMD MoE (2603.07977)
& 高------他们在 MLIP 领域影响力最大 \\
\textbf{Boris Kozinsky 组} (Harvard) & Allegro/NequIP, MoE (2604.26143)
& 高------势函数架构原创者 \\
\textbf{InstaDeep} & mlip JAX 库 (2605.22698) &
中------工程能力强，但研究侧重不同 \\
\textbf{LANL} (Sakib Matin 等) & EKD, Teacher-student RSC &
中------蒸馏方向强，但短期内不会进入多专家编译 \\
\textbf{Volker Deringer 组} (Oxford) & Foundation distillation
(2506.10956) & 中低------学术风格偏基础模型，而非工程编译 \\
\textbf{Seoul National Univ.} (Jeong Woo Han 组) & ARK Distillation &
中低------集中在催化筛选应用，而非方法论 \\
\textbf{USTC} (高玉祥等) & PFD 框架 (2502.20809) & 中------PFD
的''自动生成''方向与 EGP 互补 \\
\end{longtable}

\begin{center}\rule{0.5\linewidth}{0.5pt}\end{center}

\subsection{参考文献}\label{ux53c2ux8003ux6587ux732e}

\subsubsection{A. MoE in MLIP}\label{a.-moe-in-mlip}

\begin{enumerate}
\def\labelenumi{\arabic{enumi}.}
\tightlist
\item
  Liu, Y. et al.~``Scaling Machine Learning Interatomic Potentials with
  Mixtures of Experts.'' arXiv:2603.07977 (2026).
  https://arxiv.org/abs/2603.07977
\item
  Nascimento, G.M. et al.~``Mixture of Experts Framework in Machine
  Learning Interatomic Potentials for Atomistic Simulations.''
  arXiv:2604.26143 (2026). https://arxiv.org/abs/2604.26143
\item
  Shazeer, N. et al.~``Outrageously Large Neural Networks: The
  Sparsely-Gated Mixture-of-Experts Layer.'' ICLR 2017.
  https://arxiv.org/abs/1701.06538
\item
  Fedus, W., Zoph, B. \& Shazeer, N. ``Switch Transformers: Scaling to
  Trillion Parameter Models with Simple and Efficient Sparsity.'' JMLR
  2022. https://arxiv.org/abs/2101.03961
\item
  Birks, T., Swinburne, T. \& Kermode, J. ``Efficient and accurate
  spatial mixing of machine learned interatomic potentials.'' npj
  Computational Materials (2025).
\item
  Nascimento, G.M. \& Descoteaux, B. ``Mixture-of-Experts Framework in
  MLIP for Atomistic Simulations.'' (2026).
\end{enumerate}

\subsubsection{B. Knowledge Distillation for
MLIP}\label{b.-knowledge-distillation-for-mlip-1}

\begin{enumerate}
\def\labelenumi{\arabic{enumi}.}
\setcounter{enumi}{6}
\tightlist
\item
  ``Constructing machine learning interatomic potentials with minimum
  amount of ab initio data.'' npj Computational Materials,
  s41524-026-02023-y (2026).
  https://www.nature.com/articles/s41524-026-02023-y
\item
  Matin, S. et al.~``Ensemble Knowledge Distillation for Machine
  Learning Interatomic Potentials.'' arXiv:2503.14293 (2025).
  https://arxiv.org/abs/2503.14293
\item
  Gardner, J.L.A. et al.~``Distillation of atomistic foundation models
  across architectures and chemical domains.'' arXiv:2506.10956 (2025).
  https://arxiv.org/abs/2506.10956
\item
  Matin, S. et al.~``Teacher-student training improves the accuracy and
  efficiency of machine learning interatomic potentials.'' RSC Digital
  Discovery (2025).
  https://pubs.rsc.org/en/content/articlehtml/2025/dd/d5dd00085h
\item
  Kim, J. et al.~``A Lightweight Universal Machine-Learning Interatomic
  Potential via Knowledge Distillation for Scalable Atomistic
  Simulations (SevenNet-Nano).'' arXiv:2604.10887 (2026).
  https://arxiv.org/abs/2604.10887
\item
  Wang, R., Gao, Y. et al.~``Pre-training, fine-tuning, and distillation
  (PFD): Automatically generating machine learning force fields from
  universal models.'' Phys. Rev.~Materials 9, 113802 (2025).
  https://arxiv.org/abs/2502.20809
\item
  Lim, H. et al.~``Angular relational knowledge distillation of machine
  learning interatomic potentials for scalable catalyst exploration.''
  npj Computational Materials 12, 1 (2026).
  https://link.springer.com/article/10.1038/s41524-026-02062-5
\item
  ``Data-Efficient and Fast Machine Learning Molecular Dynamics through
  Integrated Active Learning and Knowledge Distillation.'' ChemRxiv
  (2026).
\end{enumerate}

\subsubsection{C. Expert Merging / Gauge
Alignment}\label{c.-expert-merging-gauge-alignment}

\begin{enumerate}
\def\labelenumi{\arabic{enumi}.}
\setcounter{enumi}{14}
\tightlist
\item
  ``Gauge dependence of total energies and cross-functional transfer
  learning.'' arXiv:2504.05565 (2025). https://arxiv.org/abs/2504.05565
\item
  Wortsman, M. et al.~``Model soups: averaging weights of multiple
  fine-tuned models improves accuracy without increasing inference
  time.'' ICML 2022.
\item
  Ainsworth, S. et al.~``Git Re-Basin: Merging Models modulo Permutation
  Symmetries.'' ICLR 2023.
\item
  Lee, S.H. \& Ngo, R. ``On Defining Neural Averaging.'' NeurIPS 2025.
  https://arxiv.org/abs/2508.14832
\item
  Ablin, P. et al.~``Soup-of-Experts.'' ICML 2025.
\item
  Franke et al.~``Hierarchical Re-Basin.'' ESANN 2024.
\end{enumerate}

\subsubsection{D. Governance / Arbitration (General
ML)}\label{d.-governance-arbitration-general-ml}

\begin{enumerate}
\def\labelenumi{\arabic{enumi}.}
\setcounter{enumi}{20}
\tightlist
\item
  Grassini Grimaldi, A. ``The King Solomon Protocol: A Framework for
  Human Supervised Arbitration in Multi-Model AI Systems.'' (2026).
\item
  Bie et al.~``Knowledge-guided policy arbitration: A hierarchical
  cognitive framework for safety-critical decision-making.'' Expert
  Systems with Applications (2026).
\item
  Hettiarachchi. ``From Monolithic Models to Cognitive Hives: A
  Framework for Multimodal Reasoning Systems.'' JCSTS (2026).
\item
  ``Beyond Monoliths: Expert Orchestration for More Capable, Democratic,
  and Safe Large Language Models.'' arXiv:2506.00051 (2025).
\item
  ``Toward Calibrated Mixture-of-Experts Under Distribution Shift.''
  arXiv:2606.20544 (2026).
\end{enumerate}

\begin{center}\rule{0.5\linewidth}{0.5pt}\end{center}

\begin{quote}
报告结束。本分析基于公开文献检索（截至 2026-06-25），覆盖 arXiv、Nature
Portfolio、RSC、APS、Springer 等来源。SCX/EGP
在''后训练多专家编译''这一细分领域无直接竞争，但需警惕 DeepMD 和
Kozinsky 组等团队可能的后续扩展工作。
\end{quote}

\end{document}
